\documentclass[12pt]{article}

% internal inputs
% packages
\usepackage{xcolor} % for extended colors functionality
\usepackage{parskip} % to deactivate new paragraph intendation
\usepackage{fancyhdr} % to create headers
\usepackage{graphicx} % for graphic usage
\usepackage{geometry} % for page margins
\usepackage[OT1]{fontenc} % enable european encoding and font
\usepackage{sectsty} % enable heading command formatting
\usepackage[english, ngerman]{babel} % enable syllable separation (the last language is the selected document langue)
\usepackage{colortbl} % to use colors in tables
\usepackage{tabularx} % for changing table sizes with an factor based approach
\usepackage{acro} % acronyms list, alphabetically sorted
\usepackage{pdfpages} % to import pdfs into the document
\usepackage{subcaption} % allow images arrays with subfigures
\usepackage{csquotes} % for consistent quoting
\usepackage[style=authoryear]{biblatex} % better bibliography
\usepackage{float} % allows for [H] option for figures, which actually enforces position in text
\usepackage[outputdir=../../build]{minted} % for display of code blocks
\usepackage{hyperref} % for links in general and table of content to heading linking
\usepackage{wrapfig} % enables wrap figures
\usepackage[export]{adjustbox} % enables trimming of graphics

% special characters
% http://mirrors.ctan.org/info/symbols/comprehensive/symbols-a4.pdf
\usepackage{amsfonts} % allow for \checkmark
\usepackage{pifont} % allow for \ding{228}

% dev tools
\usepackage{blindtext}

% document parameters
% you should probably never touch this file, even I forgot some parts of it
% this kinda just accumulated in configurations over the last years
% if you really need to change something... good luck...
% but i've tried to document it as good as I could!

% Colors (HEX)
% \definecolor{naklightblue}{HTML}{3cd2ff}
% \definecolor{nakblue}{HTML}{003a79}
% \definecolor{nakdarkblue}{HTML}{0f192d}
\definecolor{nakblue}{HTML}{003a79}
\definecolor{nakdarkblue}{HTML}{0f192d}
\definecolor{naklightblue}{HTML}{3cd2ff}
\definecolor{nakorange}{HTML}{ffa064}

% set graphics path
\graphicspath{{../assets/images}}

% header and footer
%% macros (need to be changed with \renewcommand)
\renewcommand{\headrulewidth}{0pt} % set thickness of a line under the header
\renewcommand{\headruleskip}{0cm} % set the distance between the line and the header text, only visual, does not impact header hight
\renewcommand{\footrulewidth}{0pt} % set thickness of a line above the footer
\renewcommand{\footruleskip}{0cm} % set the distance between the line and the footer text, only visual, does not impact footer hight

\setlength{\headsep}{1.25cm}
\setlength{\headheight}{24pt}

% page layout
\geometry{%
    a4paper,
    margin=2cm,
    top=3cm,
}

% font
\renewcommand*\familydefault{\sfdefault} % set base font of the document to sans serif

% heading formats
\newcommand{\headingStandardFormat}{\color{nakblue}\textbf{}}

\sectionfont{\headingStandardFormat}
\subsectionfont{\headingStandardFormat}
\subsubsectionfont{\headingStandardFormat}

% handle pagenumbering
%% start pagenumberin for frontmatter
\newcommand{\frontmatter}{%
    \newcounter{frontmatterPage}
    \pagenumbering{Roman} % change numbering style
}

%% start pagenumbering of mainmatter
\newcommand{\mainmatter}{%
    \setcounter{frontmatterPage}{\value{page}} % safe frontmatter counter value
    \pagenumbering{arabic} % change numbering style
}

%% reapply pagenumbering from frontmatter
\newcommand{\backmatter}{%
    \pagenumbering{Roman} % change numbering style
    \setcounter{page}{\value{frontmatterPage}} % set pagenumbering to frontmatter counter
}

\newcommand{\appendixmatter}{
    \setcounter{frontmatterPage}{\value{page}}
    \setcounter{page}{1}
    \pagenumbering{arabic}
}

\newcommand{\addAppendixTex}[3]{%
    % parameters:
    % #1: appendix heading
    % #2: filename
    % #3: label for reference
    \section{#1} % Heading for appendix
    \label{#3}
    %\addcontentsline{toc}{subsection}{#1} % register appendix in toc
    \setcounter{frontmatterPage}{\value{page} + 1} % safe page value
    \input{../assets/appendix/#2} % add pdf into main
    \backmatter{} % set correct page value, to not count appendix pages
}

% list symbol
\renewcommand{\labelitemi}{\color{nakblue}\ding{228}}

% minted

%% define highlighting style
\usemintedstyle{colorful}

%% custom listingscaption
\renewcommand{\listingscaption}{Quelltexte}

%% global minted environment style
\setminted{%
    autogobble,
    linenos, % enable line numbers
    % numbers = left, % default none
    % numbes is essentially the same as linenos exceot the side can be specified
    tabsize = 4, % default 8
    xleftmargin=3.5em
}

%% custom line numbers
\renewcommand{\theFancyVerbLine}{%
    \textcolor{gray}{%
        \ttfamily
        \arabic{FancyVerbLine}
    }
}

% hyperref

\hypersetup{%
    breaklinks=true, % default value: false
    colorlinks=true, % default value: false
    linkcolor=black, % default value: red
    % anchorcolor=black, % default value: black
    citecolor=black, % default value: green
    filecolor=black, % default value: cyan
    urlcolor=black, % default value: magenta
    bookmarksopen=true,
}

\addto\extrasngerman{%
    \crefname{section}{Abschnitt}{Abschnitte} % Singular und Plural für Anhänge
    \Crefname{section}{Abschnitt}{Abschnitte} % Großbuchstaben-Version für Anhänge
    \crefname{figure}{Abbildung}{Abbildungen}
    \Crefname{figure}{Abbildung}{Abbildungen}
    \crefname{table}{Tabelle}{Tabellen}
    \Crefname{table}{Tabelle}{Tabellen}
    \crefname{lstlisting}{Quelltext}{\listingscaption}
    \Crefname{lstlisting}{Quelltext}{\listingscaption}
    \crefname{appendix}{Anhang}{Anhänge}
    \Crefname{appendix}{Anhang}{Anhänge}
}

% tabularray

\DefTblrTemplate{contfoot-text}{default}{Fortgesetzt auf der nächsten Seite}

\NewTblrTheme{appendix}{
    \DefTblrTemplate{conthead-text}{default}{}
    \DefTblrTemplate{caption}{default}{}
    \DefTblrTemplate{conthead}{default}{}
    \DefTblrTemplate{capcont}{default}{}
}

% biblatex

\DefineBibliographyStrings{ngerman}{andothers={et al.}}
\DefineBibliographyStrings{english}{andothers={et al.}}

\renewbibmacro*{date}{%
    \iffieldundef{date}
    {\bibstring{nodate}}
    {\printdate}
}

% helper command

\newcommand{\nakQuestionTabular}[2]{%
    \renewcommand{\arraystretch}{1.5} % change vertical padding of the table

    \begin{tabularx}
        {\textwidth} % specify table width
        {| X |} % collumn rules, X fills space to collumn width
        \hline
        \rowcolor{nakblue}
        \textcolor{white}{#1} \\
        \hline
        #2                    \\
        \hline
    \end{tabularx}
}

\DeclareAcronym{BA}{%
    short = BA,
    long = Beispiel Abkürzung
}

\DeclareAcronym{ABA}{%
    short = ABA,
    long = Andere Beispiel Abkürzung
}

\def\transferleistungsNr{[NR]}
\def\matrikelnummer{[MATRIKELNUMMER]}
\def\freigegebenesThema{[THEMA]}
\def\studiengangZenturie{[STUDIENGANG, ZENTURIE]}

\begin{document}

% header
\pagestyle{fancy} % enable fancyhdr package
\fancyhf{} % clear header and footer content
\fancyhead[C]{\includegraphics[width=\textwidth]{header.png}} % add header content
\fancyfoot[C]{\thepage}

% footer
% implementing this cleanly and also looking right 
% isn't worth beeing not relevant for prettymuch everyone now
% \fancyfoot[L]{
%     \itshape{} % italic
%     * Studierende, die unter den Anwendungsbereich der PVO bis 03.02.2015 fallen, fertigen Transferleistungen
%     weiterhin in der Form von Praxisberichten an und der Begriff hält Einzug in das Abschlusszeugnis. Ab dem
%     Jahrgang 2016 hat der Begriff vollumfängliche Gültigkeit. In der Kommunikation hält der Begriff
%     Transferleistungen ab sofort Einzug.
% }

\subsection*{Selbstreflexion Transferleistung Theorie/Praxis}

\vspace{1.5em}

\renewcommand{\arraystretch}{1.5} % vertical padding

%\setlength{\tabcolsep}{0pt} % remove horozontal padding for text alignment to left margin
\begin{tabularx}{\textwidth}[h]{@{} m{.225\textwidth} @{} m{1em} X}
    Matrikelnummer:                                          &  & \matrikelnummer{}     \\
    \raggedright{} Freigegebenes Thema der Transferleistung: &  & \freigegebenesThema{} \\
\end{tabularx}

\vspace{1.5em}

Diese Fragen dienen Ihnen zur Selbstreflexion und als Unterstützung bei der Erstellung Ihrer
Transferleistung. Wir empfehlen Ihnen, sich vor, während und nach der Erstellung der
Transferleistung immer wieder diese Fragen zu stellen um auf diesem Wege eine
wissenschaftlich fundierte Transferleistung auszuarbeiten. Bitte beantworten Sie die Fragen
mit ein bis zwei Sätzen (möglich sind auch Aufzählungen) und verweisen zudem auf das
dazugehörige Kapitel in Ihrer Transferleistung. Im Gegensatz zu Ihrer Transferleistung und
Auftragsklärung ist im Rahmen der Selbstreflexion die Ich-Form erwünscht. Die Abgabe der
ausgefüllten Selbstreflexion ist Bestandteil Ihrer Transferleistung und fließt in die Bewertung
dieser ein.
\\ \\
Für die Beantwortung der Fragen stehen Ihnen als Richtwert 2000 Zeichen (ohne
Leerzeichen) zur Verfügung.

\vspace{1.5em}

\nakQuestionTabular{
    Was sind die Kernfragen meines Themas?%
}{
    \freigegebenesThema{}
}

\nakQuestionTabular{
    Welches Ziel möchte ich mit der Erstellung der Transferleistung erreichen und in welchem
    Ausmaß ist mir das am Ende gelungen?%
}{
    [ANTWORT]
}

\nakQuestionTabular{
    Wie ist die Verbindung meiner Kernfragen/meines Themas zu den im Studium erlernten Inhalten?%
}{
    [ANTWORT]
}

\nakQuestionTabular{
    Auf welche Informationen stütze ich meine Arbeit und warum habe ich diese
    Informationsbasis gewählt?%
}{
    [ANTWORT]
}

\nakQuestionTabular{
    Wie ist mein Vorgehen zur Lösung des Themas?%
}{
    [ANTWORT]
}

\nakQuestionTabular{
    Worin besteht konkret meine Eigenleistung?%
}{
    [ANTWORT]
}

\nakQuestionTabular{
    Was habe ich durch die Transferleistung gelernt und was würde ich ggf.\ beim nächsten Mal
    anders machen? Ein Verweis zu einem Kapitel ist hier optional.%
}{
    [ANTWORT]
}

\vspace{1.5em}

Laden Sie die ausgefüllte Selbstreflexion (als eigenständiges Dokument) zusammen mit Ihrer
Transferleistung im Campus Informations-System (CIS) hoch.

\end{document}
